%%%%%%%%%%%%%%%%%%%%%%%%%%%%%%%%%%%%%%%%%%%%%%%%%%%%%%%%%%%%%%%%%%%%%%%%%%%%%%%%
%%  ZWISCHENABGABE: Einlage: Eines meiner Ergebnisse   (Moodle: 20.01.2027)
%%  Aufbau nach pp-portfolio-geruest-einlage-2026-2027.tex.
%%  Der Text steht in einlagen/4_messergebnis.tex und wird 1:1 in die Endabgabe übernommen.
%%  Overleaf: Menü -> Main document = diese Datei. PDF für Moodle benennen:
%%  Manevski_Messergebnis.pdf
%%%%%%%%%%%%%%%%%%%%%%%%%%%%%%%%%%%%%%%%%%%%%%%%%%%%%%%%%%%%%%%%%%%%%%%%%%%%%%%%
\newcommand{\dasDatum}{20.\,01.\,2027}   % Moodle-Frist dieser Abgabe
%%%%%%%%%%%%%%%%%%%%%%%%%%%%%%%%%%%%%%%%%%%%%%%%%%%%%%%%%%%%%%%%%%%%%%%%%%%%%%%%
%%  EINSTELLUNGEN – die einzige Datei, in der du deine Angaben einträgst.
%%  Wird von ALLEN Hauptdateien (Endabgabe + Zwischenabgaben) eingelesen.
%%  Entspricht dem Block "Die Angaben" der offiziellen Gerüste
%%  pp-portfolio-geruest-2026-2027.tex / pp-portfolio-geruest-einlage-2026-2027.tex
%%%%%%%%%%%%%%%%%%%%%%%%%%%%%%%%%%%%%%%%%%%%%%%%%%%%%%%%%%%%%%%%%%%%%%%%%%%%%%%%

%%%%%%%%%%%%%%%%%%%%%%%%%%%%%%%%%%%%% Die Angaben %%%%%%%%%%%%%%%%%%%%%%%%%%%%%%%
\newcommand{\Semester}{WH 2026-2027}
\newcommand{\Projekt}{Dynamik Keltischer Wackelsteine}   % Titel laut Themenliste WS 26/27
%% Schwerpunkt abgeleitet aus der Aufgabenverteilung im Arbeitsplan (1.4 Bewegungsgleichungen,
%% 1.6 Simulation Umkehrdynamik, 2.1.3/2.1.4 Firmware + Kalibrierung, 5.1 Datenpipeline) – prüfen!
\newcommand{\Schwerpunkt}{Bewegungsgleichungen und numerische Simulation der Spin-Umkehr}
\newcommand{\SchwerpunktKurz}{Simulation der Spin-Umkehr}        % kurze Form für die Kopfzeile
\newcommand{\IhrName}{Kristijan Manevski}                         % Vorname Nachname
\newcommand{\Matrikelnummer}{\platzhalter{Matrikelnummer}}
\newcommand{\Studiengang}{\platzhalter{Studiengang + Semester}}   % z.B. Physik B.Sc., 5. Semester
\newcommand{\emailAdresse}{\platzhalter{Uni-E-Mail-Adresse (\dots @uni-ulm.de)}}
%% Abgabedatum: jede Zwischenabgabe-Hauptdatei setzt ihr eigenes \dasDatum (Moodle-Frist);
%% ohne Vorgabe gilt das Datum der Endabgabe.
\providecommand{\dasDatum}{15.\,03.\,2027}

%% Weitere Angaben, die in mehreren Einlagen gebraucht werden
\newcommand{\Partner}{Janosch Lustig}                             % für Abgrenzung der Schwerpunkte
%% laut Titelseite von Janoschs Arbeitsplan-Abgabe (28.10.2026)
\newcommand{\PartnerSchwerpunkt}{Umkehrzeit und Parameterabhängigkeit}
\newcommand{\Laborbuchlink}{https://link-zum-elektronischen-laborjournal}  % Pflicht in Einlage 3!
%% Einlage 3: ein Praktikumstag (Do 10--12 Uhr, O26/233 bzw. O26/234), ≠ Tag von Janosch.
%% Termine: 15.10., 22.10., 29.10., 05.11., 12.11., 19.11., 26.11., 03.12., 10.12., 17.12.2026,
%%          14.01., 21.01., 28.01., 04.02., 11.02.2027 (Abgabe Protokoll: 09.12.2026)
\newcommand{\Protokolldatum}{TT.MM.2026}
\newcommand{\PaperKey}{KondoNakanishi2017}      % Einlage 2: Kondo & Nakanishi, Phys. Rev. E 95, 062207
\newcommand{\Vortragsdatum}{28.\,01.\,2027}      % Einlage 5
\newcommand{\FremdVortragKey}{VORLAGE-Vortrag}  % Einlage 6: bib-Key des anderen Vortrags

%%%%%%%%%%%%%%%%%%%%%%%%%%%%%%%%%%%%%%%%%%%%%%%%%%%%%%%%%%%%%%%%%%%%%%%%%%%%%%%%
%%  MODUS-SCHALTER
%%
%%  \Abgabefalse  = Entwurf: grüne Hilfen sichtbar, Kontrollseite mit
%%                  Umfangsmessung, Platzhalter rot markiert.
%%  \Abgabetrue   = Abgabe: grüne Farbe -> schwarz (wie offizielle Zeile 15/16),
%%                  alle Hilfen + Kontrollseite + Umfangsboxen verschwinden.
%%                  Übrig gebliebene \platzhalter{} bleiben ROT sichtbar und
%%                  erzeugen eine Warnung -> dann ist noch etwas offen!
%%
%%  \Hinweisefalse blendet im Entwurf nur die grünen Hilfstexte aus
%%                  (nützlich, um den Umfang realistisch zu messen).
%%%%%%%%%%%%%%%%%%%%%%%%%%%%%%%%%%%%%%%%%%%%%%%%%%%%%%%%%%%%%%%%%%%%%%%%%%%%%%%%
\newif\ifAbgabe
\newif\ifHinweise

\Abgabefalse      % <-- für die Abgabe auf \Abgabetrue setzen
\Hinweisetrue     % <-- Hilfstexte im Entwurf ein-/ausblenden

\newcommand{\dieEinlage}{Einlage: Eines meiner Ergebnisse}

\documentclass[pdftex,12pt,twoside,footsepline,headsepline,a4paper,final]{scrreprt}
%%%%%%%%%%%%%%%%%%%%%%%%%%%%%%%%%%%%%%%%%%%%%%%%%%%%%%%%%%%%%%%%%%%%%%%%%%%%%%%%
%%  PRÄAMBEL – gemeinsam für Endabgabe und alle Zwischenabgaben.
%%  Teil 1 ist 1:1 die Präambel der offiziellen Gerüste (nicht verändern),
%%  Teil 2 ergänzt nur Pakete für Gleichungen, Bilder, PDF-Einbindung usw.,
%%  Teil 3 enthält die Hilfsmakros dieser Vorlage.
%%  Die \documentclass-Zeile steht (unverändert offiziell) in jeder Hauptdatei.
%%%%%%%%%%%%%%%%%%%%%%%%%%%%%%%%%%%%%%%%%%%%%%%%%%%%%%%%%%%%%%%%%%%%%%%%%%%%%%%%

%%%%%%%%%%%%%%%%%%%%%%%%%%%%% Teil 1: offizielle Präambel %%%%%%%%%%%%%%%%%%%%%%
\usepackage[verbose,a4paper,bottom=25mm,top=25mm]{geometry} %showframe,
\usepackage{scrlayer-scrpage}
\usepackage[T1]{fontenc}
\usepackage[utf8]{inputenc}
\usepackage{ngerman}
\usepackage{color}
\usepackage{amsfonts}
\usepackage{amssymb}
\usepackage{amsthm}
%%%%%%%%%%%%%%%%%%%%%%%%% Biblatex (offizielle Optionen, Backend Biber) %%%%%%%%%%
%% Overleaf erkennt biber automatisch. Zitierstil alphabetic: [Sch02]
\usepackage[sorting=nyt, backend=biber, citestyle=alphabetic, style=alphabetic,
            autocite=inline, labelalpha=true, hyperref=auto, backref=true,
            language=ngerman]{biblatex}
\addbibresource{literatur.bib}

%%%%%%%%%%%%%%%%%%%%%%%%%%%%% Teil 2: Zusatzpakete %%%%%%%%%%%%%%%%%%%%%%%%%%%%%%
\usepackage{amsmath}               % \eqref, align, nummerierte Gleichungen
\usepackage{xcolor}
\definecolor{darkgreen}{rgb}{0,0.5,0}   % offizielle Hinweisfarbe
\usepackage{graphicx}              % \includegraphics
\graphicspath{{abbildungen/}}
\usepackage{booktabs}              % schöne Tabellen
\usepackage[locale=DE]{siunitx}    % \SI{3.4}{\kilo\hertz}, \num{1.2e3}
\usepackage{pdfpages}              % \includepdf für Anhänge (Paper, Laborjournal, Folien)
\usepackage{pdflscape}             % Querformat: \begin{landscape}...\end{landscape}
\usepackage{pgfgantt}              % Gantt-Diagramm für den Arbeitsplan
\usepackage{environ}
\usepackage[savepos,abspage]{zref}
\usepackage[hidelinks]{hyperref}   % klickbare Verweise; zuletzt laden

%% Zitierregel (Marti): bei Journalartikeln DOI angeben, aber keine Webadresse.
%% -> url/urldate werden bei @article automatisch unterdrückt.
\AtEveryBibitem{%
  \ifentrytype{article}{\clearfield{url}\clearfield{urlyear}%
    \clearfield{urlmonth}\clearfield{urlday}}{}}

%% Kopf- und Fußzeile wie offizielles Gerüst (\ohead wird in der Hauptdatei gesetzt)
\pagestyle{scrheadings}
\ihead{\thepage}
\ofoot{\pagemark}
\ifoot{\zubearbeiten{\IhrName}}

%%%%%%%%%%%%%%%%%%%%%%%%%%%%% Teil 3: Hilfsmakros %%%%%%%%%%%%%%%%%%%%%%%%%%%%%%%
\makeatletter

%% --- offizieller Farbschalter (Zeile 12-15 der Gerüste), automatisch per Modus
\ifAbgabe
  \Hinweisefalse
  \newcommand{\zubearbeiten}[1]{\textcolor{black}{#1}}
\else
  \newcommand{\zubearbeiten}[1]{\textcolor{darkgreen}{#1}}
\fi

%% --- Grüne Hilfstexte: verschwinden im Abgabemodus vollständig
%%     \begin{hinweise} ... \end{hinweise}   (darf itemize enthalten)
%%     \hinweis{kurzer Text}
\NewEnviron{hinweise}{\ifHinweise\par{\color{darkgreen}\small\BODY\par}\fi}
\newcommand{\hinweis}[1]{\ifHinweise\par{\color{darkgreen}\small #1\par}\fi}

%% --- Platzhalter für eigenen Text: bleibt IMMER rot sichtbar (auch im
%%     Abgabemodus), damit nichts Unfertiges abgegeben wird.
\newcounter{platzhalter}
\DeclareRobustCommand{\platzhalter}[1]{%
  \stepcounter{platzhalter}%
  \textcolor{red}{\textbf{[}\,#1\,\textbf{]}}%
  \ifAbgabe\PackageWarning{Portfolio}{Platzhalter noch vorhanden: #1}\fi}
\AtEndDocument{%
  \ifnum\value{platzhalter}>0
    \PackageWarningNoLine{Portfolio}{Es sind noch \theplatzhalter\space
      Platzhalter im Dokument!}%
  \fi
  \immediate\write\@auxout{\string\gdef\string\pf@platzhalteralt{\theplatzhalter}}}
\providecommand{\pf@platzhalteralt}{?}

%% --- Umfangsmessung
%%   \begin{umfang}{Bezeichnung}{min}{max} ... \end{umfang}
%%   misst die Länge des eingeschlossenen Textes in Seiten (über die
%%   Position auf der Seite, auf 0.05 Seiten genau) und zeigt sie im Entwurf
%%   in einem Kasten an. Achtung: Abbildungen/Tabellen zählen nach den Regeln
%%   NICHT zur Länge, werden hier aber mitgemessen -> dann abziehen.
\zref@addprop{savepos}{abspage}
\newcounter{pf@umfang}
\newcommand*{\pf@umfangliste}{}
\newcommand*{\pf@umfangeintrag}[4]{% Name, min, max, ist
  \g@addto@macro\pf@umfangliste{\pf@umfangzeile{#1}{#2}{#3}{#4}}}
\newcommand*{\pf@setfarbe}[3]{% min max ist -> \pf@farbe
  \def\pf@farbe{darkgreen}%
  \ifdim #3pt<#1pt \def\pf@farbe{red}\fi
  \ifdim #3pt>#2pt \def\pf@farbe{red}\fi}
\newcommand*{\pf@umfangzeile}[4]{%
  #1 & #2 -- #3 & \pf@setfarbe{#2}{#3}{#4}\textcolor{\pf@farbe}{\textbf{#4}}\\}
\newenvironment{umfang}[3]{%
  \stepcounter{pf@umfang}%
  \edef\pf@uid{pfumf\the\value{pf@umfang}}%
  \def\pf@uname{#1}\def\pf@umin{#2}\def\pf@umax{#3}%
  \par\zsavepos{\pf@uid a}\nobreak
}{%
  \par\zsavepos{\pf@uid e}%
  \ifAbgabe\else
    \zref@refused{\pf@uid a}\zref@refused{\pf@uid e}%
    \edef\pf@ist{\fpeval{round(
        (\zref@extractdefault{\pf@uid e}{abspage}{0}
        -\zref@extractdefault{\pf@uid a}{abspage}{0})
        +(\zref@extractdefault{\pf@uid a}{posy}{0}
        -\zref@extractdefault{\pf@uid e}{posy}{0})/\number\textheight ,2)}}%
    \protected@write\@auxout{}{\string\pf@umfangeintrag{\pf@uname}{\pf@umin}{\pf@umax}{\pf@ist}}%
    \par\smallskip\noindent
    \pf@setfarbe{\pf@umin}{\pf@umax}{\pf@ist}%
    \fcolorbox{\pf@farbe}{white}{\scriptsize
      Umfang \emph{\pf@uname}: $\approx$\,\pf@ist~Seiten
      (Soll: \pf@umin\,--\,\pf@umax~Seiten, ohne Titel/Abb./Tab./Legenden)}%
    \par\smallskip
  \fi}

%% --- Kontrollseite (nur im Entwurf; ersetzt das grüne Kapitel
%%     "Hinweise zur Arbeit mit dieser Datei", das abgegeben nicht vorkommen darf)
\newcommand{\kontrollseite}{%
  \ifAbgabe\else
  \chapter*{Kontrollseite (erscheint nur im Entwurf)}
  {\color{darkgreen}
  Modus: \textbf{Entwurf}. Für die Abgabe in \texttt{einstellungen.tex}
  \texttt{\textbackslash Abgabetrue} setzen. Dann wird alles Grüne schwarz bzw.\
  entfernt (Regel: \enquote{Abgaben mit grüner Farbe oder stehen gebliebenen
  Hilfen ergeben einen Abzug}).\par\medskip
  \textbf{Offene Platzhalter} (letzter Durchlauf): \pf@platzhalteralt\par\medskip
  \textbf{Umfangskontrolle} (Werte aus dem letzten Kompilierdurchlauf;
  Hilfstexte stehen außerhalb der Messung):\par\smallskip
  \ifx\pf@umfangliste\@empty
    \emph{noch keine Messwerte -- nochmals kompilieren.}
  \else
    {\small\begin{tabular}{p{0.55\textwidth}cc}
      \toprule Teil & Soll [Seiten] & Ist [Seiten]\\\midrule
      \pf@umfangliste
      \bottomrule
    \end{tabular}}
  \fi\par\medskip
  %% Checkliste – erscheint nur auf der Kontrollseite im Entwurf.
%% Quellen: Portfolios im Projektpraktikum 2026-2027 (Marti, 09.07.2026),
%% offizielle Gerüste 2026-2027, Zitierregeln (31.08.2026),
%% Abbildungen und Gleichungen (Limmer 2007), Moodle-Terminliste.

\textbf{Fristen (Moodle; nächste Abgabe wird erst nach der vorherigen freigeschaltet!)}\par\smallskip
{\small\begin{tabular}{@{}ll@{}}
21.\,10.\,2026 & Motivationsformular (ohne Reflexion) -- obligatorisch\\
28.\,10.\,2026 & Arbeitsplan + Reflexion -- obligatorisch\\
11.\,11.\,2026 & Wissenschaftliche Arbeit + Reflexion -- obligatorisch\\
09.\,12.\,2026 & Protokoll eines Praktikumstages + Reflexion -- obligatorisch\\
20.\,01.\,2027 & Messergebnis + Reflexion -- obligatorisch\\
23.\,01.\,2027 & Späteste Anmeldung mündliche Leistung (Portal)\\
27.\,01.\,2027 & Folien des Vortrags + Materialien für Mitstudierende\\
28.\,01.\,2027 & Vorträge (Prüfung, Anwesenheitspflicht)\\
07.\,02.\,2027 & Folien mit Kommentar + Reflexion (vorgeschlagen, freiwillig)\\
21.\,02.\,2027 & Aussage aus anderem Vortrag + Reflexion (vorgeschlagen, freiwillig)\\
10.\,03.\,2027 & Späteste Anmeldung schriftliche Leistung (Portal)\\
\textbf{15.\,03.\,2027} & \textbf{Endabgabe Portfolio -- verspätet = nicht bestanden}\\
\end{tabular}}\par\smallskip
{\small Praktikumstage (Google Kalender): donnerstags 10--12\,Uhr, O26/233 bzw.\ O26/234 --
15.\,10.\ bis 17.\,12.\,2026 und 14.\,01.\ bis 11.\,02.\,2027; keine Termine 18.\,12.--13.\,01.
Die Anmeldetermine im Hochschulportal (23.\,01./10.\,03.) stammen aus dem Plan von Petros -- im Portal prüfen!}
\par\medskip

\textbf{Vor jeder Abgabe prüfen}
\begin{itemize}\itemsep0pt
\item \texttt{\textbackslash Abgabetrue} gesetzt, keine grüne Farbe, keine Hilfen, \emph{0 Platzhalter}.
\item Deckblatt vollständig: Name, Vorname, Matrikelnummer, Studiengang, E-Mail, Thema (Projekt + Schwerpunkt).
\item Inhaltsverzeichnis mit Seitenzahlen, auch der Anhänge; durchgehend Seitenzahlen.
\item Längen eingehalten (siehe Tabelle oben; Abbildungen/Tabellen/Legenden zählen nicht mit):
      Einleitung 2--4, Beschreibung $\tfrac14$--$\tfrac12$, Diskussion der Physik 1--3,
      Reflexion 1--2, Gesamtreflexion 3--6 Seiten. Schriftgröße 12\,pt (fest eingestellt).
\item Beschreibung und Reflexion sind \emph{keine Nacherzählung}; die Leitfragen werden in Fließtext
      beantwortet (nicht in Einzelsätzen), erkennbar zugeordnet.
\item Paper (Einlage 2) wird in den späteren Einlagen zur Argumentation verwendet; Gleichungen werden
      zur Argumentation genutzt.
\item Peer-Feedback (Kommilitonen) eingeholt, aufgegriffen -- und zitiert. Petros' Arbeitsplan als Vorlage
      zitieren (Key \texttt{Petros2026Arbeitsplan}); Schwerpunkt-Abgrenzung mit Janosch festhalten.
\item Jede Übernahme zitiert: Literatur, Bilder (\texttt{\textbackslash cite[Abb.~3.4]\{key\}}),
      Gleichungen (\texttt{\textbackslash cite[Gl.~(3)]\{key\}}), Aussagen von Personen, Anleitungen,
      Datenblätter, Webseiten (mit Abrufdatum), \textbf{jede KI-Nutzung mit System und ALLEN Prompts}.
\item Journalartikel: Autoren, Titel, Journal, Band, Seiten, Jahr, DOI (nicht als Link).
\item Abbildungen: Bildunterschrift, im Text referenziert (\texttt{\textbackslash ref}), Schrift
      $\geq 10$\,pt, Achsen mit Einheit als \enquote{$f$ / kHz} (nicht \enquote{[kHz]}), ggf.\ Residuenplot.
\item Tabellen: Überschrift, im Text referenziert; Messdaten nicht als lange Listen im Text, sondern im Anhang.
\item Wichtige Gleichungen nummeriert, alle Symbole definiert und einheitlich verwendet.
\item Dateiname für Moodle: \texttt{Manevski\_WelcheEinlage.pdf} (nur PDF oder DOCX).
\item Bis auf den Arbeitsplan müssen sich alle Einlagen von denen der Partnerin/des Partners unterscheiden.
\end{itemize}

  }%
  \fi}
\providecommand{\enquote}[1]{\glqq #1\grqq}

%% --- Anhang aus PDF mit durchgehenden Seitenzahlen (Regel: "Die Arbeit ist
%%     durchgängig mit Seitenzahlen versehen ... auch der Anhänge")
%%   \anhangpdf{Kapiteltitel}{label}{datei.pdf}           alle Seiten
%%   \anhangpdf[pages={1-3}]{Kapiteltitel}{label}{datei.pdf}
%%   Die Kapitelüberschrift steht auf einer eigenen Seite, danach folgen die
%%   PDF-Seiten (auf 82 % verkleinert, damit Kopf-/Fußzeile frei bleiben).
\newcommand{\anhangpdf}[4][pages=-]{%
  \chapter{#2}\label{#3}%
  \includepdf[#1,scale=0.82,pagecommand={\thispagestyle{scrheadings}}]{#4}}
%% Gleiches ohne neues Kapitel (für weitere PDFs im selben Anhang)
\newcommand{\pdfseiten}[2][pages=-]{%
  \includepdf[#1,scale=0.82,pagecommand={\thispagestyle{scrheadings}}]{#2}}

\makeatother

%% Antwortfeld im Motivationsformular (Anhang der Endabgabe)
\newcommand{\antwort}[1]{\par\smallskip #1\par\bigskip}

\ohead{\zubearbeiten{\Projekt{} - \SchwerpunktKurz}}

\begin{document}

%% Titelseite einer ZWISCHENABGABE – wörtlich aus pp-portfolio-geruest-einlage-2026-2027.tex
%% \dieEinlage wird in der jeweiligen Hauptdatei (Manevski_....tex) gesetzt.
\title{Universit\"{a}t Ulm\\\vspace{1cm}
{\LARGE \textnormal{\textsc{Fakult\"{a}t f\"{u}r Naturwissenschaften}}}\\\vspace{2cm} {\LARGE \textnormal{\textsl{Projektpraktikum \Semester}}}\\\vspace{2cm} {\LARGE\zubearbeiten{\textit{\Projekt}}}\\
{\Large \zubearbeiten{\textit{\Schwerpunkt}}}\\\vspace{0.5cm} \zubearbeiten{\textit{\dieEinlage}}}

\author{{\LARGE \zubearbeiten{\IhrName}}\\ $_{}$\\$_{}$\\\vspace*{1.cm}
Matrikelnummer:  \zubearbeiten{\Matrikelnummer}\\\vspace{0.8cm} Studiengang: \zubearbeiten{\Studiengang}\\\vspace{0.4cm} e-Mail: \zubearbeiten{\emailAdresse}}

\date{\zubearbeiten{\dasDatum}}

\maketitle

\tableofcontents
\kontrollseite   % nur im Entwurf sichtbar

\chapter{Hauptteil}
\begin{hinweise}
Beachten Sie die Hinweise in der Datei mit dem Gerüst für die Endabgabe! Auch bei einer Zwischenabgabe ist
eine korrekte Literaturarbeit notwendig. Denken Sie daran, dass Sie alle Hinweise im Text auf andere Werke
(Artikel, Bücher, Bilder, Grafiken) aber auch alle Erwähnungen von Meinungen anderer Personen mit einem Zitat
versehen müssen! Beachten Sie auch die Fragen bei der jeweiligen Abgabe in Moodle.
\end{hinweise}

%%%%%%%%%%%%%%%%%%%%%%%%%%%%%%%%%%%%%%%%%%%%%%%%%%%%%%%%%%%%%%%%%%%%%%%%%%%%%%%%
%%  EINLAGE 4: EINES IHRER ERGEBNISSE     Zwischenabgabe: 20.01.2027 (obligatorisch)
%%  Struktur: Beschreibung (1/4–1/2 S.) + Diskussion der Physik (1–3 S.) + Reflexion (1–2 S.)
%%  Rohdaten, Auswerteformeln, ausgewertete Daten, Grafiken gehören zur Einlage.
%%  Lange Messwertlisten -> Anhang (4_messergebnis_anhang.tex), nicht in den Text.
%%%%%%%%%%%%%%%%%%%%%%%%%%%%%%%%%%%%%%%%%%%%%%%%%%%%%%%%%%%%%%%%%%%%%%%%%%%%%%%%
\section{Einlage: Eines Ihrer Ergebnisse (Messung oder Berechnung).}
\label{sec:messergebnis}

\begin{hinweise}
Hier nehmen Sie ein von Ihnen gemessenes Resultat (Rohdaten, Auswerteformeln und die ausgewerteten Daten und
Grafiken) und diskutieren es. Ein Bezug zu Ihrem Paper (Einlage~2) ist Voraussetzung.
\end{hinweise}

\subsection{Beschreibung}
\begin{hinweise}
Die Beschreibung gibt alle notwendigen Informationen an, die die Einlage beschreiben, also zum Beispiel das
Datum, die Art der Messung und eine kurze Inhaltsübersicht, aber keine Nacherzählung (min. $1/4$ Seite bis
max. $1/2$ Seite!).
\end{hinweise}
\begin{umfang}{Einlage 4 (Messergebnis): Beschreibung}{0.25}{0.5}
\platzhalter{Datum, Art der Messung, Messaufbau in Kürze -- laut Arbeitsplan: 3D-gedruckter Wackelstein
mit eingebautem Inertialsensor (Gyroskop) auf nivellierter Glasplatte (30\,cm\,$\times$\,30\,cm),
reproduzierbarer Anstoß mit definiertem Anfangs-$\omega$; eine der Messreihen 1--3 (Anfangs-$\omega$
variiert, Nick- vs.\ Roll-Anregung, $\delta$ variiert); Auswertung mit der eigenen Python-Datenpipeline;
Rohdaten in Anhang~\ref{anh:messdaten}, Ergebnis in Abbildung~\ref{fig:messergebnis}}
\end{umfang}

\subsection{Diskussion der Physik}
\begin{hinweise}
Beschreiben Sie hier die Physik dieser Einlage.
\begin{itemize}\itemsep0pt
\item Was haben Sie gemessen oder berechnet?
\item Was ist die Hauptaussage?
\item Was sind die Voraussetzungen für Ihre Messung?
\item Welches sind Ihre Auswerteformeln?
\item Was trägt ihre \textbf{wissenschaftliche Arbeit} zur Argumentation bei?
\item Welche Physik dieser Einlage ist bei Ihrem Schwerpunktthema anwendbar?
\item Wo gab es Schwierigkeiten beim Verständnis der Physik? Welche?
\item Wie haben Sie Ihr Paper aus der Einlage \flqq Eine wissenschaftliche Arbeit\frqq{} für diese Einlage verwendet? Warum? Wo?
\item Stimmen Ihre Ergebnisse mit den Aussagen aus dem Paper überein?
\end{itemize}
\textbf{Die Diskussion der Physik sollte mindestens eine Seite und maximal 3 Seiten lang sein. Die
Längenangaben sind ohne Titel, Bilder, Tabellen und Tabellen- oder Bildlegenden und ohne übergrosse Abstände
zu Gleichungen zu verstehen. Ich werde gegebenenfalls Ihre Auswertung anhand der beigelegten Daten
überprüfen.}

\smallskip Regeln für Abbildungen/Tabellen/Gleichungen (Limmer): jede Abbildung mit Bildunterschrift und im
Text referenziert; Schrift in Diagrammen $\geq 10$\,pt in Endgröße; Achsen \enquote{Größe / Einheit}
(z.\,B. $\omega_z$\,/\,rad\,s$^{-1}$, nicht $\omega_z$ [rad/s]); bei guter Übereinstimmung
Messung--Theorie/Simulation zusätzlich Residuenplot; Tabellen mit Überschrift; Messdaten nicht (nur) als
Tabelle, sondern (auch) als Diagramm; wichtige Gleichungen nummeriert, alle Symbole definiert.
Die Gerüste unten (Gleichung, Abbildung, Tabelle) sind auf das Projekt zugeschnitten -- anpassen oder löschen.
Bezug zum Paper ist Pflicht: Messergebnis mit der Umkehrzeit-Formel bzw.\ den Simulationen aus
\cite{\PaperKey} und mit der eigenen Simulation (Arbeitspaket 1.6) vergleichen.
\end{hinweise}
\begin{umfang}{Einlage 4 (Messergebnis): Diskussion der Physik}{1}{3}
\platzhalter{Was wurde gemessen, Hauptaussage, Voraussetzungen (Rollen ohne Gleiten auf der Glasplatte,
Kalibrierung des Gyroskops, \dots)}

% ---------- Gerüst: Auswerteformel mit Nummer und Verweis ----------
% Wachstum der Nick- bzw. Rollschwingung in der instabilen Drehrichtung (Arbeitspaket 5.2):
\begin{equation}
  A(t) = A_0\,\mathrm{e}^{\gamma t},
  \label{eq:wachstum}
\end{equation}
wobei $A$ die Amplitude der Nick- bzw.\ Rollschwingung, $A_0$ ihr Anfangswert und $\gamma$ die
Wachstumsrate der Instabilität ist. \platzhalter{Gleichung prüfen/ersetzen; Umkehrzeit $t_r$ aus den
Gyroskopdaten definieren (z.\,B.\ Zeit bis zum Vorzeichenwechsel von $\omega_z$); Bezug zum Paper:
\enquote{Nach \cite[Gl.~(\dots)]{\PaperKey} gilt für $t_r$ \dots}; Gl.~\eqref{eq:wachstum} im Text verwenden}

% ---------- Gerüst: Abbildung (Datei nach abbildungen/ legen) ----------
\begin{figure}[htbp]
  \centering
  %\includegraphics[width=0.8\textwidth]{messergebnis.pdf}
  \fbox{\parbox[c][4cm][c]{0.7\textwidth}{\centering\platzhalter{Abbildung: abbildungen/messergebnis.pdf,
    z.\,B.\ $\omega_z(t)$ aus dem Gyroskop mit Spin-Umkehr}}}
  \caption{\platzhalter{Aussagekräftige Bildunterschrift: Winkelgeschwindigkeit $\omega_z$ um die
    Hochachse über der Zeit $t$ für \dots, Messbedingungen, Ergebnis der Anpassung mit
    Gl.~\eqref{eq:wachstum}. Fremde Abbildungen: \cite[Abb.~\dots]{\PaperKey}}}
  \label{fig:messergebnis}
\end{figure}

% ---------- Gerüst: kurze Ergebnistabelle (Rohdaten in den Anhang!) ----------
\begin{table}[htbp]
  \centering
  \caption{\platzhalter{Tabellenüberschrift, z.\,B.\ gemessene Umkehrzeit $t_r$ in Abhängigkeit
    vom Anfangs-$\omega_0$ (Messreihe 1) im Vergleich zu Simulation und Paper}}
  \label{tab:ergebnis}
  \begin{tabular}{cccc}
    \toprule
    {$\omega_0$ / \si{\radian\per\second}} & {$t_r$ (Messung) / \si{\second}}
      & {$t_r$ (Simulation) / \si{\second}} & {$t_r$ (Paper) / \si{\second}} \\
    \midrule
    \dots & \dots & \dots & \dots \\
    \bottomrule
  \end{tabular}
\end{table}

\platzhalter{Diskussion: Abbildung~\ref{fig:messergebnis} und Tabelle~\ref{tab:ergebnis} im Text
verwenden, Unsicherheiten (Fehlerrechnung, Arbeitspaket 5.3), Vergleich Messung -- eigene Simulation --
Paper (5.4), Anwendbarkeit auf \Schwerpunkt}
\end{umfang}

\subsection{Reflexion}
\begin{hinweise}
In der Reflexion sollen Sie sich Gedanken machen über die Wirkung dieser Einlage und der Arbeit mit ihr auf
Ihren Lernfortschritt.
\begin{itemize}\itemsep0pt
\item In welchem Zusammenhang steht die Einlage mit Ihrem Schwerpunktthema?
\item Was haben Sie bei der Erstellung dieser Einlage gelernt?
\item Wie hat Ihnen die Einlage beim Verständnis Ihres Projektthemas geholfen?
\item Wo gab es Schwierigkeiten bei der Beschäftigung mit der Einlage?
\item Was hätten Sie anders gemacht, wenn Sie die Einlage nochmals erstellen müssten?
\item Was sagen Ihre Kommilitoninnen oder Kommilitonen dazu, als Kommentar oder wenn Sie mit Ihnen über die Einlage diskutieren?
\item Haben Sie in der Durchführung oder Auswertung der Messung logische oder physikalische Fehler entdeckt?
\item Gab es bei der Beschäftigung mit der Einlage eine überraschende neue Erkenntnis?
\end{itemize}
\textbf{Die Reflexion für eine Einlage sollte mindestens eine Seite und maximal 2 Seiten lang sein. Die
Längenangaben sind ohne Titel, Bilder, Tabellen und Tabellen- oder Bildlegenden zu verstehen.}
\end{hinweise}
\begin{umfang}{Einlage 4 (Messergebnis): Reflexion}{1}{2}
\platzhalter{Reflexion in Fließtext -- Leitfragen erkennbar beantworten}
\end{umfang}

%\subsection{Eigene Bewertungskriterien}
%...


\printbibliography[heading=bibintoc,title={Referenzen und Literaturverzeichnis}]

\appendix
\chapter*{Anh\"{a}nge}
\addcontentsline{toc}{chapter}{Anh\"{a}nge}
%%  ANHANG zu Einlage 4: Rohdaten und Auswertung
%%  ("Ich werde gegebenenfalls Ihre Auswertung anhand der beigelegten Daten überprüfen.")
%%  Möglichkeiten: Tabelle(n) hier, PDF-Export (\pdfseiten), Auswerteskript als Listing.
%%  Große Datenmengen zusätzlich als Datei in Moodle bzw. Link zum Laborjournal.

\chapter{Messdaten und Auswertung (Einlage 4)}
\label{anh:messdaten}

\platzhalter{Rohdaten des Inertialsensors (Gyroskop, ggf.\ Beschleunigung) der gewählten Messreihe,
Kalibrierwerte (Gyro-Bias, Achsen), Körperparameter des Messkörpers (Masse, Schwerpunkt, Trägheitstensor),
Geräte/Einstellungen (Abtastrate, Glasplatte), Python-Datenpipeline bzw.\ Auswerteschritte}

%\pdfseiten{anhaenge/messdaten.pdf}


\end{document}
